\documentclass[11pt,a4paper]{article}
\usepackage[utf8]{inputenc}
\usepackage[T1]{fontenc}
\usepackage{amsmath,amssymb,amsthm}
\usepackage{geometry}
\geometry{margin=2.5cm}
\theoremstyle{plain}
\newtheorem{theorem}{Theorem}[section]
\newtheorem{proposition}[theorem]{Proposition}
\newtheorem{lemma}[theorem]{Lemma}
\theoremstyle{definition}
\newtheorem{definition}[theorem]{Definition}
\title{Soluciones Tests 57-61: Ecuaciones Diferenciales}
\author{Mathematical AI Agent}
\date{\today}
\begin{document}
\maketitle
\section{Problemas}

\subsection{Test 57: Ecuación diferencial exponencial}
\begin{theorem}[Solución de $y' = ky$]
Sea $k \in \mathbb{R}$ y $y_0 \in \mathbb{R}$. La única solución del problema de valor inicial
\[
y' = ky, \quad y(0) = y_0
\]
es $y(t) = y_0 e^{kt}$ para todo $t \in \mathbb{R}$.
\end{theorem}

\begin{proof}
Definimos $u(t) = y(t) e^{-kt}$. Derivando con respecto a $t$:
\[
u'(t) = y'(t) e^{-kt} - k\, y(t) e^{-kt} = (y'(t) - k\, y(t)) e^{-kt}.
\]
Como $y$ satisface $y' = ky$, se tiene que $y'(t) - ky(t) = 0$ para todo $t$, luego:
\[
u'(t) = 0 \cdot e^{-kt} = 0.
\]
Por lo tanto, $u$ es constante: $u(t) = u(0)$ para todo $t$. Evaluando en $t=0$:
\[
u(0) = y(0) e^{0} = y_0.
\]
Así, $y(t) e^{-kt} = y_0$, de donde $y(t) = y_0 e^{kt}$.

Para la unicidad, supongamos que $z$ es otra solución. Entonces $w(t) = y(t) - z(t)$ satisface $w' = kw$ con $w(0) = 0$. Por el mismo argumento anterior, $w(t) = w(0) e^{kt} = 0$ para todo $t$, luego $y = z$.
\end{proof}

\subsection{Test 58: EDO afín}
\begin{theorem}[Solución de $y' = ay + b$]
Sea $a, b \in \mathbb{R}$ con $a \neq 0$. La solución general de la ecuación diferencial
\[
y' = ay + b
\]
es
\[
y(t) = C e^{at} - \frac{b}{a}
\]
donde $C \in \mathbb{R}$ es una constante arbitraria.
\end{theorem}

\begin{proof}
Definimos la función constante $y_p(t) = -b/a$. Se verifica directamente que:
\[
y_p'(t) = 0 = a \left(-\frac{b}{a}\right) + b = -b + b = 0.
\]
Por lo tanto, $y_p$ es una solución particular de la ecuación.

Sea $y$ cualquier solución. Definimos $u(t) = y(t) - y_p(t) = y(t) + b/a$. Entonces:
\[
u'(t) = y'(t) = a\, y(t) + b = a\left(u(t) - \frac{b}{a}\right) + b = a\, u(t).
\]
Así, $u$ satisface $u' = au$, cuya solución general por el Test 57 es $u(t) = C e^{at}$. Por lo tanto:
\[
y(t) = u(t) + y_p(t) = C e^{at} - \frac{b}{a}.
\]
\end{proof}

\subsection{Test 59: Unicidad con condición de Lipschitz}
\begin{theorem}[Unicidad bajo condición de Lipschitz]
Sea $f: D \subseteq \mathbb{R} \times \mathbb{R}^n \to \mathbb{R}^n$ continua en un dominio $D$ y supongamos que existe una constante $L > 0$ tal que para todo $(t, y_1), (t, y_2) \in D$:
\[
\|f(t, y_1) - f(t, y_2)\| \leq L \|y_1 - y_2\|.
\]
Entonces el problema de valor inicial
\[
y' = f(t, y), \quad y(t_0) = y_0
\]
tiene exactamente una solución local en algún intervalo $[t_0 - h, t_0 + h]$.
\end{theorem}

\begin{proof}
\textbf{Existencia (Picard--Lindelöf):} Definimos la sucesión de iteradas de Picard:
\[
y_0(t) = y_0, \quad y_{n+1}(t) = y_0 + \int_{t_0}^{t} f(s, y_n(s))\, ds.
\]
Sea $M = \max_{|t - t_0| \leq h} \|f(t, y_0)\|$ y elegimos $h > 0$ tal que $hL < 1$. Probamos por inducción que para $n \geq 1$:
\[
\|y_{n+1}(t) - y_n(t)\| \leq \frac{(Lh)^n}{n!} M h.
\]
El caso base ($n=0$) se verifica por definición. Para el paso inductivo:
\begin{align*}
\|y_{n+2}(t) - y_{n+1}(t)\| &\leq \int_{t_0}^{t} \|f(s, y_{n+1}(s)) - f(s, y_n(s))\|\, ds \\
&\leq L \int_{t_0}^{t} \|y_{n+1}(s) - y_n(s)\|\, ds \\
&\leq L \int_{t_0}^{t} \frac{(Lh)^n}{n!} M h\, ds = \frac{(Lh)^{n+1}}{(n+1)!} M h.
\end{align*}
La serie $\sum_{n=0}^{\infty} \frac{(Lh)^n}{n!} M h$ converge, por lo que la sucesión $(y_n)$ converge uniformemente a una función $y$ continua. Pasando al límite en la integral:
\[
y(t) = y_0 + \int_{t_0}^{t} f(s, y(s))\, ds,
\]
y por el Teorema Fundamental del Cálculo, $y$ es solución de la EDO con $y(t_0) = y_0$.

\textbf{Unicidad:} Sean $y$ y $z$ dos soluciones. Sea $w(t) = y(t) - z(t)$. Entonces $w(t_0) = 0$ y para $t \geq t_0$:
\[
\|w(t)\| = \left\|\int_{t_0}^{t} (f(s, y(s)) - f(s, z(s)))\, ds\right\| \leq L \int_{t_0}^{t} \|w(s)\|\, ds.
\]
Por la desigualdad de Gronwall generalizada, $\|w(t)\| \leq 0$ para $t \geq t_0$, luego $w(t) = 0$. Análogamente para $t \leq t_0$, obteniendo $y = z$.
\end{proof}

\subsection{Test 60: Sistemas lineales y matriz exponencial}
\begin{theorem}[Solución mediante matriz exponencial]
Sea $A \in \mathbb{R}^{n \times n}$ una matriz constante. Toda solución del sistema lineal
\[
x' = Ax
\]
se puede expresar como
\[
x(t) = e^{At} x_0
\]
donde $x_0 = x(0)$ y $e^{At} = \sum_{k=0}^{\infty} \frac{(At)^k}{k!}$ es la matriz exponencial.
\end{theorem}

\begin{proof}
\textbf{Paso 1: $e^{At}$ es una matriz de Lie.} Probamos que $\frac{d}{dt} e^{At} = A e^{At}$. Por definición de la serie de potencias:
\[
\frac{d}{dt} e^{At} = \frac{d}{dt} \sum_{k=0}^{\infty} \frac{A^k t^k}{k!} = \sum_{k=1}^{\infty} \frac{A^k t^{k-1}}{(k-1)!} = A \sum_{j=0}^{\infty} \frac{A^j t^j}{j!} = A e^{At}.
\]
La convergencia de la serie permite la derivación término a término.

\textbf{Paso 2: $e^{At}$ es la solución fundamental.} Definimos $\Phi(t) = e^{At}$. Entonces $\Phi'(t) = A \Phi(t)$ y $\Phi(0) = I$ (la matriz identidad). Para cualquier $x_0 \in \mathbb{R}^n$, definimos $x(t) = \Phi(t) x_0 = e^{At} x_0$. Entonces:
\[
x'(t) = \Phi'(t) x_0 = A \Phi(t) x_0 = A x(t),
\]
y $x(0) = \Phi(0) x_0 = x_0$. Así, $x$ es solución del problema de valor inicial.

\textbf{Paso 3: Unicidad.} Sea $y(t)$ cualquier solución de $x' = Ax$ con $y(0) = x_0$. Definimos $z(t) = e^{-At} y(t)$. Entonces:
\begin{align*}
z'(t) &= -A e^{-At} y(t) + e^{-At} y'(t) = -A e^{-At} y(t) + e^{-At} A y(t) \\
&= e^{-At}(-A y(t) + A y(t)) = 0,
\end{align*}
donde se usó que $A$ y $e^{-At}$ conmutan. Por lo tanto, $z(t) = z(0) = y(0) = x_0$, luego $y(t) = e^{At} x_0$.
\end{proof}

\subsection{Test 61: Estabilidad asintótica}
\begin{theorem}[Estabilidad asintótica por autovalores]
Sea $A \in \mathbb{R}^{n \times n}$ y supongamos que todos los autovalores $\lambda_1, \ldots, \lambda_n$ de $A$ satisfacen $\operatorname{Re}(\lambda_i) < 0$ para todo $i = 1, \ldots, n$. Entonces el sistema lineal $x' = Ax$ es asintóticamente estable, es decir, toda solución $x(t)$ satisface $\lim_{t \to \infty} x(t) = 0$.
\end{theorem}

\begin{proof}
Sea $x(t) = e^{At} x_0$ la solución del sistema con condición inicial $x_0$. Debemos demostrar que $\|e^{At}\| \to 0$ cuando $t \to \infty$, donde $\|\cdot\|$ denota una norma matricial inducida.

\textbf{Caso 1: $A$ es diagonalizable.} Entonces existe una matriz invertible $P$ tal que $A = P D P^{-1}$ con $D = \operatorname{diag}(\lambda_1, \ldots, \lambda_n)$. Entonces:
\[
e^{At} = P e^{Dt} P^{-1} = P \operatorname{diag}(e^{\lambda_1 t}, \ldots, e^{\lambda_n t}) P^{-1}.
\]
Para cada autovalor $\lambda_i = \alpha_i + i \beta_i$ con $\alpha_i < 0$:
\[
|e^{\lambda_i t}| = e^{\alpha_i t} \to 0 \quad \text{cuando } t \to \infty.
\]
Por lo tanto, $\|e^{Dt}\| = \max_i |e^{\lambda_i t}| \to 0$, y como:
\[
\|e^{At}\| \leq \|P\| \cdot \|e^{Dt}\| \cdot \|P^{-1}\|,
\]
se concluye que $\|e^{At}\| \to 0$.

\textbf{Caso 2: $A$ general.} Para toda $\varepsilon > 0$, existe una constante $C_\varepsilon > 0$ tal que para todo $t \geq 0$:
\[
\|e^{At}\| \leq C_\varepsilon e^{(\mu + \varepsilon) t}
\]
donde $\mu = \max_i \operatorname{Re}(\lambda_i) < 0$. Elegimos $\varepsilon > 0$ suficientemente pequeño tal que $\mu + \varepsilon < 0$. Entonces:
\[
\|e^{At}\| \leq C_\varepsilon e^{(\mu + \varepsilon) t} \to 0 \quad \text{cuando } t \to \infty.
\]
En particular, para toda solución $x(t) = e^{At} x_0$:
\[
\|x(t)\| \leq \|e^{At}\| \cdot \|x_0\| \to 0,
\]
lo que prueba la estabilidad asintótica.
\end{proof}

\end{document}
